% !TeX root = ../incremental_SS_Translation.tex

\section{Uttaratantra 64: Preservation of Health}% svastharakṣanīya

\section{Literature}

Meulenbeld offered an annotated overview of this chapter, which is called
\emph{svasthavṛtta} in the vulgate.\fvolcite{IA}[330--331]{meul-hist} There,
he translates \emph{svasthavṛtta} as ``preservation of health'' but elsewhere
as ``maintenance of health''\fvolcite{IA}[301]{meul-hist} and
``hygiene.''\fvolcite{IB}[738]{meul-hist} His only reference to secondary
literature on this topic is to an annotated German translation by J. Laping
1984.\footnote{\cite{lapi-1984}, cited by \volcite{IB}[431]{meul-hist}.} His
footnotes to the chapter also indicate
many parallels with \emph{Carakasaṃhitā, Sūtrasthāna}, chapter six.

On the term \emph{svastha}, Meulenbeld cites W. Halbfass (1991): 250-252; M.
Hara (1995) and P.V. Sharma (1985e: 52-53).\fvolcite{IB}[10]{meul-hist}

In primary sources, see the definition of \emph{svasthā} in
\emph{Suśrutasaṃhitā, Sūtrasthāna} 15.41--44; \emph{svasthahita} in
\emph{Aṣṭāṅgahṛdaya, Sūtrasthāna} 1.16cd; ...

Comparable chapters in other medical texts include the \emph{svasthacatuṣka}
in \emph{Carakasaṃhitā, Sūtrasthāna} (see below) and a chapter on
\emph{svasthavṛtti} in Cakrapāṇidatta's \emph{Cakradatta}.

Meulenbeld also refers to \emph{svasthacatuṣka}, a group of four chapters
(5--8) in the \emph{Carakasaṃhitā, Sūtrasthāna}, chapters 5–8, also called
\emph{svasthavṛttacatuṣka} or \emph{svāsthyavṛttikacatuṣka} in various
colophons of the text and commentaries. He summarises these
chapters,\fvolcite{IA}[13--14]{meul-hist} and their names are
\emph{mātrāśitīya}, daily regime (cf. Bhela Sū.8; A.h.Sū.8; A.s.Sū.11. Cf.
Su.Ci.24); \emph{tasyāśitīya}, seasonal regimen (\emph{ṛtucaryā})
(Meulenbeld cites many references in primary and secondary
literature\fvolcite{IB}[13]{meul-hist}); \emph{navegāndhāraṇīya}, the
unwholesome effects of suppressing natural urges, etc. (Cf. Bhela Sū.6;
Su.Ci.24); and \emph{indriyopakramaṇīya}, senses, mind and rules of conduct
(Cf. Bhela Sū.7; Su.Ci.24. See on Ca.Sū.8: S.Ch. Vidyabhusana (1971): 11--12).

Meulenbeld mentions that a ``Sadānanda wrote a commentary, called \emph{Auṣadhavivṛti}, on the chapters of the
\emph{Carakasaṃhitā} containing rules for the preservation of health
(svasthavṛtta).\fvolcite{IA}[196]{meul-hist} His footnote on this work refers to
the NCC (VI, 396–397) and the following bibliographic reference:
\emph{pañcatantram \ldots\
	carakasūtrasthānasvasthavṛttacatuṣkākhyacaturadhyāyāḥ\ldots\
	sadānandaśāstrikṛtauṣadhavivṛtiyutayā saṃvalitaṃ}, Mercantile Press, Lahore
	1926 [IO.San.D.554].\fvolcite{IB}[302]{meul-hist}

In the vulgate, this chapter is included in a section called, ``the chapters that
are an embellishment of the text" (\emph{tantrabhūṣaṇādhyāya}). However, in
the Nepalese version, it is the eighth part of tenfold section on treatment for the
body (\emph{kāyacikitsita}).\footnote{See witness K, folio 209r, ll. 4--5
	(\dev{mūtradoṣo mūtrāghāto yonyamānūṣa eva ca/
		apasmāronmādakañcaiva rasabhedastathaiva ca//
		svastharakṣāvidhāṇañca tantrayuktiśca doṣabhit/
		ityebhirdaśakaṃ proktaṃ punaḥ kāyacikitsite//}).}

This chapter is cast as a discussion between Dhanvantarī and the king, who is not named but is likely to be Divodāsa.

% History of Health Preservation in Indian Medicine 

% Outline

% What is Health (Premodern South Asia)?
% * Definitions of svasthā, svasthya and svastha (the healthy person) in premodern Sanskrit literature.
% * svastha vs rogyatā (Aṣṭāṅgahṛdaya 1.20). An essentialist claim. 
% * What are the Normative assumptions?
% * How does one define health in medicine more broadly? On the Normal and the Pathological, Georges Canguilhem, 1991. The Nature of Disease, Lawrie Reznek, 1987. Disease in History, William McNeil (article). Ivan Illich, Medical Nemesis, 1975 (medicine has become a threat to health).
% * Are there characteristics of health (or diseases) in Ayurveda that are not considered as health in other cultures?
% * does svasthavṛtti allow one to fulfil one's life's goal? Yes, Cakrapāṇi on Caraka links them (See Dominik's article on Medicine and Dharma).
% * Individual vs. social health
% * Lived experience? Literature that records people's views on svastha and their healthy regimes.
% * Is it age related? Does it cover mental health (unmada and apasmāra)?

% How was health preserved?
% a. Seasonal living (do people live seasonally in 21st century?)
% * vidhis (dos and don'ts)
% b. Dietary regimes
% c. Living space and clothing
% d. Physical exertion (exercise and sex). Edit and translate all textual passages on vyāyāma.  

% Health Prevention vs Treatment
% The focus of svastharakṣanīya is clearly the prevention of disease, but to what extent are the same elements of its regime (diet, exercise, etc.) incorporated into the treatment of sick people? Is there much of an overlap?

% Comparative Section?
% a. Concept of svastha vs:
% 	Chinese
% 	Greek/Roman
% 	Byzantine/Persian?

% 1. Primary Sources
% a. Suśruta: Nepalese Version vs Vulgate
% b. Caraka (Sūtrasthāna, the svasthacatuṣka, chs. 5–8)
% c. Vāgbhaṭa
% d. Cakrapāṇi's Cakradatta (chapter on svasthavṛtti)
% e. Commentators: Ḍalhaṇa (with relevant citations of Jejaṭa and Gayadāsa), Cakrapāṇi on Caraka, Sadānanda's Auṣadhavivṛti

% 2. Edited Chapters of the Nepalese Version of the Suśrutasaṃhitā
% a. Uttaratantra 64 (svastharakṣanīya)
% b. Sūtrasthāna 6 (ṛtucaryā): Edit and translate the whole chapter, but with a view to discussing in the introduction selected sections on definitions of the seasons, their characteristics; seasonal variations of doṣas; diseases caused by seasonal doṣas. In some ways, this chapter should precede 6.64.
% c. Sūtrasthāna 46 (annapānavidhi):  Edit and translation only the section on seasonal diet.


\section{Translation}
% K f. 205v
% H f. 406

\begin{translation}

	\item [1]
	Now we shall explain how to preserve a person's health.%
	% the upkeep of a healthy person.% rakṣaṇīya - maintenance of? prophylaxis for a healthy person?
	\footnote{The term \emph{svastharakṣaṇīya} more literally means, ``preserve a healthy person.'' \Dalhana{6.64.1}{808} noted this reading, which is also what we see in the Nepalese version, and explained it in these terms.

		The word \emph{rakṣaṇīya} (preservation, maintenance, upkeep, etc.) is replaced in the vulgate \citep[808]{vulgate} by \emph{vṛtta}, which, when compounded with \emph{svastha}, denotes the “regime for a healthy person.” However, \emph{rakṣaṇa} is retained later in the same chapter of the vulgate.}

	\item [2]
	One whose humours, digestive fire, and the functioning of bodily constituents and waste products are in equilibrium; whose self, senses, and mind are serene is called healthy.%
	\footnote{This verse is absent here in the vulgate but occurs at the end of \emph{Sūtrasthāna} 15 in both the Nepalese version and vulgate. The vulgate \citep[808]{vulgate} supplements its omission with the statement, ``Such a healthy person has been described in the \emph{Sūtrasthāna}'' (\emph{sūtrasthāne samuddiṣṭaḥ svastho bhavati yādṛśaḥ}). \Dalhana{1.15.41}{75--76} discussed several points of view regarding the meaning of equilibrium in these different contexts.}\q{More on Ḍalhaṇa and the meaning of sama?}

	\item [3]
	My dear, the preservation of a healthy person is the supreme benefit of medicine.\footnote{%
		Subsequent verses in this chapter indicate that this is a discussion between Dhanvantari and a king, who is likely to be Divodāsa.} %
	The topic of his regime and preservation, which was mentioned by me at the beginning, was outlined in brief and will here be explained in detail.
	% The interlocutors are king (divodāsa) and Dhanvantari.See the end of the chapter.
	% something has gone wrong with artham in the Nepalese manuscripts (which has arthāṃ and arthān). The participle (uktaṃ) and verb (vakṣyate) require a singluar. And artham can be neuter. The syntax of the vulgate is wrong because the subject is plural (arthāḥ) and the verb singular (a patch made inadequate by metric causa?)
	% ḍalhaṇa: āditaḥ anāgatabādhe

	\item [4]
	In each season, when specific humours are aggravated in people, a wise physician should prescribe the appropriate tastes for that season.

	\subsubsection{Wet Weather Regimen}
	\item [5] % varṣa season
	You should know that when people's digestive fire is weak because of the dampness of their bodies during the rainy months, wind and other humours become aggravated due to their excitement.\footnote{%
		The vulgate notes that the readings \emph{saṃgharṣāt} and \emph{saṃharṣāt} occur in an unspecified number of witnesses (\cite[809, n.\,2]{vulgate}). The reading \emph{saṃgharṣāt} (``friction'') makes better sense in the context of the humours, but is not supported by the Nepalese manuscripts.}
	% Ḍalhaṇa appears to understand mārutādayaḥ as referring to the doṣas.
	% have I over-translated khalu?

	\item [6--8ab]
	Therefore, in order to remove moisture and subdue the humours,
	the king should be given solid food with astringent, bitter, and \se{kaṭu}{spicy} tastes, which is neither too oily nor too dry, and which is hot and \se
	{dīpana}{stimulates digestion}.\footnote{%
		The vulgate reads \emph{doṣasaṃ\textbf{hara}ṇāya} (``for removing the humours'') instead of \emph{doṣasaṃ\textbf{dhāra}ṇāya} \citep[809]{vulgate}. \CS\ \Ca{2.6.6}{219--220} uses the word \emph{saṃdhāraṇa} in the context of restraining the urges to urinate and so on. In the current context, \emph{doṣasaṃdhāraṇa} refers to subduing the humours that have been aggravated by weak digestive fire during the rainy season (mentioned in the previous verse).}
	He should drink the water which was mentioned before,\footnote{%
		\Dalhana{6.64.8}{809} understood “the water which was mentioned before” to be monsoon rainwater, which was described in the section on what to drink (\dev{yajjalaṃ coktamādita iti ‘varṣāsvantarikṣaṃ’ ityādinā pānīyavarge kathitam}). This is the section that starts at \SS\ \Su{1.45.3–46}{196--200}.} %
	or water which has been treated with hot iron\q{%
		Look for references to treating water with hot iron.} %
	and mixed with honey.%
	%

	% See Roots Ed 3 p 228 for the tastes and the passage in the Aṣṭāṅgahṛdaya on the tastes
	% 
	% Seasons in North India: there are six:
	% dewy season, %śiśira 
	% spring, % vasanta
	% summer, % grīṣma 
	% the monsoon season, % varṣā 
	% autumn, % śarad 
	% and winter. % hemanta

	\item [8cd--9]
	On a rainy and windy day that is very cold and wet, % atyarthaśītāmbusaṃkule: literally: affected by excessively cold water
	herbal medicines develop acidity because of being fresh.%
	\footnote{Ḍalhaṇa \citep[809]{vulgate}  glossed \emph{vidāham} (“burning”) as \emph{amlapāka} (“sour digestion”), suggesting that fresh herbs may cause acidic irritation in the stomach when taken on a rainy, windy, cold, damp day.}
	%
	And because of that, a sensible person should not exercise excessively.

	\item [10] % not in the vulgate
	When it is not very cold, aged barley, wheat and rice with spicy broths and soups, as well as appetizing meats from dry land,
	should be eaten.\footnote{``The meats of dry land'' are those of \sepl{jaṅgala}{arid land}; see \cite[fig.\,5, p.\,52]{zimm-1999} for a list of the \SS's meats in this class.\label{jaṅgala}}
	% emended to bhojyā based on Caraka Sū 6.38c (note in the edition) 
	% the ca in pāda 4 suggests that jāṅgala is wild meat/deer/game rather than an adjective. Elsewhere in Ḍalhaṇa I see eṇādayo jāṅgalamṛgamāṃseṣu (an adjective here but obviously refers to game).

	\item [11] % not in the vulgate
	One should drink a little mead or sweet wine,%
	\footnote{On \emph{mādhvīka}, a kind of honey liquour, see McHugh 2021: 46. On \emph{sīdhu} (“sweet wine”), Ḍalhaṇa \citep[778]{vulgate} said it is \emph{gauḍika}, an alcoholic drink produced from sugar or grain (\emph{gauḍikaṃ sīdhu guḍakṛtaṃ sasyakam}). McHugh (2021: 45–46), who translates \emph{sīdhu} as “sugarcane wine,” says that it is primarily made from sugarcane but also other sweet substances like jaggery, herbs and mahua flowers.}
	%
	and also rainwater, lake water, or well water, heated and cooled and mixed with honey.\footnote{The same idea in similar words is expressed in \CS\ \Ca{1.6.38--39}{47}.}

	\item[12] % ab not in the vulgate
	At this time, % atra or 'In this climate'?
	one should avoid the easterly wind, exposure to  heat, \se{udamantha}{buttered barley water}, river water, sleep during the day, frost, and sex.

	\item[13--14ab] % not in the vulgate
	One should have \se{pragharṣa}{body rubs}, \se{udvartana}{massages}, baths, perfumes, and garlands and one should smoke. To avoid humidity from the ground, one should have a heavy cloak and sleep in a house's central room, which is cool, has a fire, and is free from draughts.
	%
	% uṣma heat, hot vapour, steam. Is "humidity" a better translation  here as I think this still concerns the varṣa season?
	% cf: Caraka1.6.40: pragharṣodvartanasnānagandhamālyaparo bhavet | laghuśuddhāmbaraḥ sthānaṃ bhajed akledi vārṣikam || [during the rainy season]
	% cf: AS.Sū.4.48	pragharṣodvartanasnānadhūmagandhāgarupriyaḥ | 
	%				yāyāt kareṇumukhyābhiścitrasragvastrabhūṣitaḥ || [varṣa season]
	% Jayaratha on Tantrāloka: madhyegṛhamiti--gṛhamadhye brahmasthāne--ity arthaḥ
	% madhyamaveśmani might denote 'in an inner/central chamber of the house.'
	% pragharṣa = gharṣaṇa (Āyurvedīya Kośa), and gharṣaṇa = mardana. mardana = aṅgamardana
	% udvartanaṃ = gātramardanaṃ (Āyurvedīya Kośa); 1. tvagavagharṣaṇa 2. snānāt pūrvaṃ tvagavagharṣaṇārtham upayujyamānaṃ cūrṇam āmalakyādi. 3. pādābhyām aṅgāvayavānāṃ pīḍanaṃ yena vyāyāmajanitaḥ śramaḥ śāmyati.

	\item[14c–f] % not in the vulgate
	And if one is injured, one should avoid \sepl{doṣa}{morbid conditions}, such as damp, and be protected by an umbrella. One also should travel by means of female elephants and adorn oneself with high-quality \gls{aguru}.%
	\footnote{Ḍalhaṇa \citep[809]{vulgate} remarks that one should travel on female elephants to avoid mud and water, and one should wear aloe-wood because it wards off the cold (\emph{nāgavadhūbhir iti hastinīgamanaṃ ca kardamāmbunoḥ parihārārtham| praśastāgurubhūṣita iti aguruṇaḥ, śaityanivārakatvāt}).} %
	So it is.%
	\footnote{The \emph{iti} after 15ab appears to follow from verse 3. In K, the \emph{iti} is followed by a rubricated section marker.}\\

	%
	% Then the vulgate has:
	% divāsvapnam ajīrṇaṃ ca varjayet tatra yatnataḥ | 
	% sevyāḥ śaradi yatnena kaṣāyasvādutiktakāḥ ||
	% kṣīrekṣuvikṛtikṣaudraśālimudgādijāṅgalāḥ |
	% śvetasrajaś candrapādāḥ pradoṣe laghu cāmbaram ||
	% salilaṃ ca prasannatvāt sarvam eva tadā hitam
	% saraḥsv āplavanaṃ caiva kamalotpalaśāliṣu ||

	\subsubsection{Autumn} % ?? śarad ?? The beginning of autumn is marked by the end of the monsoon followed by relatively cooler weather.

	\item[15--16]
	Homage to the sun % bhāskare namaḥ
	which reaches its extent
	% labdhaprasara: attested in MW: mf(\skt{ā})n. that which has obtained free scope, moving at liberty, unimpeded Mudr. But does it  refer to the sun rising (i.e., when the sun has attained its rising) or the spreading forth of its light?
	along the bright pathways of the quarters, % prasannadiṅmārga
	when the sky covered % avaguṇṭhite vyomni literally: when the sky is enveloped by ghananīla: % not sure about °vaguṇṭhite vyomni. I've translated avaguṇṭhite more like a noun. 
	with clouds and dark blue % ghananīla: ghana = clouds, nīla = dark blue.
	across the expanse of the earth has dispersed. % vyapanītadharābhoga: bhoga or ābhoga here as curve or expanse: the expanse of the earth
	% On dharābhoga: cf. sasvarge sadharābhoge sapātāle sadiktaṭe / jagaty asmiṃs tu sarvasmin na kiñcic chobhanaṃ sthiram // MU_6,24.29
	% vyapanīta means removes or eliminates
	It is exalted by being born at the sight of Indra’s thunderbolt. % I have not found a reference to a śakrāstra, and its relation to the sun.
	% I am assuming the locative is being used for the dative. Oberlies observes that namaskṛtya can take the locative in Epics (p 353), but also that the locative and dative are exchangeable. 
	The splendour of the sun advances, smothering cloud after cloud.

	\item[17]
	Bile that is accumulated during the rains is aggravated by the heat in autumn. For this reason, when all the clouds have dispersed one should undertake practices that eliminate bile.

	\item[18]

	During autumn, when one has an appetite, one should eat food that is cool, astringent and sweet, slightly light and  bitter, as well as mixed with ghee and containing meat of animals from arid regions.%
	\footnote{On the meaning of \emph{jāṅgala} (``the meat of animals from arid regions''), see footnote \ref{jaṅgala}.}
	% can food be both astringent and sweet or should we understand the compound kaṣāyamadhuram as astringent or sweet?
	% I have read  īṣal in compound with laghu. 
	% Is satiktakam bitter or pungent?
	% jāṅgala appears to range from deer (jaṅghāla) to birds and burrowing animals (Ḍalhaṇa 1.46.53: jāṅgalaśabdena jaṅghālaviṣkirapratudaparṇamṛgabileśayā iti trayodaśa bhedāḥ ṣaṭsvevāntarbhūtāḥ)

	\item[19]
	While not indulging excessively in sex, one should consume \gls{hariṇa}, % deer
	\gls{pṛṣata}, %chital, 
	\glspl{lāva}, % quails, 
	\glspl{chaga}, %  goats, 
	\glspl{eṇa}, % black antelopes, 
	\gls{kapiñjala} % partridges,
	\glspl{mudga}, % mung beans, 
	\gls{śāli}, and aged \gls{yava}. % barley.
	% Are there two ways to read this? One having lots of sex should not consume.
	% And should jīrṇān be read with mudgāñ chālīn and yavāñ, or only yavāñ 

	\item[20]
	%\begin{sloka}
	One should avoid burning, acrid, spicy and heating things, sleeping during the day, and heavy foods, as well as vegetables and meats that are heavy, and exposure to frost.
	%\end{sloka}

	% It seems strange to have divāśvapna in the middle of a compound of types of food, so I suspect that at least one anusvāra has dropped out, more likely two (unless I've misunderstood gurūṇi. Note my emendations. The atha suggests at least two compounds were intended here (instead of the one in K and H).
	% Cf. Aṣṭāṅgahṛdaya, nidānasthāna 13.2526
	% vyādhikarmopavāsādikṣīṇasya bhajato drutam | 
	% atimātram athānyasya gurvamlasnigdhaśītalam | 
	% lavaṇakṣāratīkṣṇoṣṇaśākāmbu svapnajāgaram |
	% mṛdgrāmyamāṃsavallūram ajīrṇaśramamaithunam ||

	\item[21] % not in the vulgate
	In autumn, O King, the wise person should carefully follow this rule:
	at that time one should consume grapes, sugar cane and milk. %(\emph{tatra}).
	% confirmation that a king is being spoken to here.

	\item[22] % xml:id 6.64.21
	When bile has accumulated during the rains, one should eliminate it with purgatives, or by applications of bitter medicated ghees, and also by bloodletting.
	% I assume that tiktaḥ and sarpis need to be read in compound with prayoga (and have emended tikta accordingly). 

	\item[23] % xml:id = 6.64.18cd
	% In the vulgate marked by round brackets
	Water should be tasty and cool, pure, as clear as a flawless crystal, cleansed by the rays of the autumn sun, and detoxified by the rising star Agastya. Because of its clearness, all such water is then beneficial.
	% See Vitus Angermeier's PhD Das Wasser thesis on how water is purified. 1880 Holtzmann article on Agastya ZDMG



	\item[24] % xml:id 6.64.19a-21
	%  In the vulgate marked by round brackets
	And in autumn seasons, one should enjoy the breeze amongst lotuses, lilies and rice plants or sandalwood with fragrant camphor; and clean, light clothes;%
	\footnote{The vulgate reading of the text is, ``one should enjoy clean, light clothing [perfumed] with sandalwood and camphor'' \citep[809]{vulgate}. There are various references in early Sanskrit literature to sandalwood powder being used to fragrance clothes (\emph{paṭavāsa(ka)}), but this is not the case for camphor, the possible exception being a concoction for perfuming clothes in \emph{Bṛhatsaṃhitā} 76.12 (\volcite{2}[841]{trip-1968}) that has camphor as an ingredient.} %
	autumnal garlands; and drinking sweet wine judiciously.
	%
	%
	% pavanam ... kamalotpalaśāliṣu: the breeze blowing among lotuses, lilies and rice stalks?
	% yuktitaḥ – within reason? reasonably? sensibly

	\item[25] % 6.64.21ab
	% 26cd not in vulgate
	One should do everything that pacifies bile: enjoying the moonlight, a circle of friends, pleasant comforts, and sweet speech. Thus, [autumn].%\footnote{Verses in the vulgate not in the Nepalese version.}
	\q{fill out this note.}

	\subsubsection{Winter} % % 6.64.21cd

	\item[26--27]
	Winter is cold and dry. Its sun is mild and the wind is blustery. %
	One should consume the following food and drinks:
	\gls{tila} seeds,
	\glspl{māṣa},
	vegetables and curds, as well as
	sugarcane products, and
	rices that are fragrant and also fresh, %
	%
	and meats % māṃsāni
	of carnivorous % prasaha
	and marsh-dwelling animals, % anūpa or ānūpa
	meats of
	creatures that are
	carnivorous and live in burrows, % kravyādabilaśāyinām
	aquatic animals and % audakānām
	birds, % plava = MW: m. a kind of aquatic bird
	and amphibious creatures.%
	\footnote{The syntax of verse 27 is rather odd. It appears that the author intended the reader to supply \dev{māṃsāni} with the genitive plurals in 27b–d. The possibility that the genitives qualify the compound \dev{prasahānūpamāṃsāni} is excluded by the fact that a list of animals in  \SS\ \Su{1.46.135--136}{225--226}  includes \dev{prasaha}, \dev{anūpa}, \dev{kravyāda}, and \dev{bilavāsa} separately, indicating that they are to be understood as different types of animals.} % pādināñ 
	% Emended māṣāṃ śākāni (K) to māṣāñ śākāni
	% emended vikṛtīñ (K: pc) to vikṛtīḥ (vulgate).
	% I don't understand the syntax here as the genitive plurals appear to depend on māṃsāni, but māṃsāni is the final member of a compound. Maybe the original reading was something like prasahānāṃ ca māṃsāni.
	% I've translated kravyāda as an adjective. Not sure if that's right, but we also have prasaha as carnivorous animals.

	\item[28]
	If one wishes to fortify oneself at the onset of winter, %puṣṭim icchan himāgame
	one may consume as desired % niṣeveta...kāmatas
	clear alcoholic drinks, % madyāni ca prasannāni
	and anything that gives strength. %  yac ca kiñcid balapradam

	\item[29] % 31cd and 32ab are absent in the vulgate
	Bathing % avagāhanam
	in moderately or comfortably warm  % taile koṣṇe sukhoṣṇe vā
	oil is recommended. % praśastam
	Furthermore, % tathā
	one should engage in \se{abhyaṅga}{oil massage}, including the head, and perform vigorous exercise.  Alternatively one should  bathe in tepid water and wash with only warm water.
	% vā is bit awkward

	\item[30--31]
	One whose body has been smeared with \gls{kuṅkuma} and \gls{aguru} % kuṃkumāgurudigdhāṃgo. Is this referring to some sort of paste made with saffron and aloes-wood?
	should rest % śayīta
	in the large %mahati
	inner chamber of a house,  % gṛhagahvare
	that is free from draughts % nivāte
	and covered with a carpet. % kuthakāstṛte. kutha: Apte - carpet (general)
	It should have a cart of coals, % sāṅgārayāne
	be filled with incense, perfume, and musk,  % dhūpāmodamadotkaṭau: mada; musk. āmoda; strong fragrance. utkaṭa; richly endowed with, abounding in. 
	and have a bed covered by wool and silk fabrics.% aurṇṇakauśeyasamvītaśayane
	\footnote{%
		Ḍalhaṇa \citep[810]{vulgate} glosses \dev{sāṅgārayāna} (``vessel for coals'') as \dev{aṅgārapūrṇaśakaṭikāsahita} (``with a small cart filled with coals''). On \dev{kauśeya}, Ḍalhaṇa \citep[86]{vulgate} says it is produced by the saliva of a worm that makes a cocoon (\dev{kośakārakṛmilālodbhūta}).
	} %
	And one should rest on a bed that is % cāpi śayane. 
	well covered % suvistīrṇṇe could mean "well-covered" (with blankets, etc.). However, vistīrṇa would normally mean wide.
	and beautiful. % manorame

	\item[32] % not in the vulgate
	There, % tatra
	affectionate % priyā
	women who are well-adorned, % nāryaḥ svalaṃkṛtāḥ
	have cast aside their necklaces, % apanītahārās: is hāra more specifically a pearl necklace?
	and intensely passionate, % madotkaṭāḥ
	should delight him, % rāmayeyur
	at the proper time, % yathākālaṃ
	even by force, % balād api
	while he is being attended to %sevyamānaṃ
	with soft touches % mṛduśparśair
	and insistent embraces. % nirdayair upagūhanaiḥ: nirdaya - merciless, cruel, violent, excessive
	% Cakrapāṇi quotes part of this verse: suśrute+api hemante maithunasevoktā; yaduktaṃ- “tatrāpanītahārāśca priyā nāryaḥ svalaṅkṛtāḥ| ramayeyur yathākāmaṃ balād api madotkaṭāḥ” iti||9-18|| This justifies two of the emendations I have made.


	% got to here DW & JB
	\subsubsection{Late Winter}
	% remaining part of winter: the dewy season (śiśira). Also translated as "cold season", "late winter". Maybe hemanta = early winter, śiśira = late winter

	\item[33]
	This regime % eṣa vidhi
	that has been taught should be done during late winter (\emph{śiśira}). % śiśira: the cool or dewy season (comprising two months, Māgha and Phālguna, or from about the middle of January to that of March)
	% Is the point of this statement that thhis vidhi should also be done during śiśira (and not just hemanta) or that it should especially be done during śiśira.
	Some physicians consider \emph{śiśira} to be the remaining part of winter.

	\item[34] % not in the vulgate
	Owing to the cold [of late winter] produced by clouds and wind and the dry caused by the warm seasons,%
	\footnote{The term \emph{ādāna} has been translated as ``warm seasons'' because it is generally understood to be the warm half of the year or the three warm seasons during which the sun takes (ā + ⁄{dā}) strength from living beings. The term appears in \emph{Carakasaṃhitā, Sūtrasthāna} 6.4, where it is defined as the three seasons beginning with late winter (\emph{śiśira}) and ending with summer (\emph{griṣma}). The  \emph{Āyurevedīyakośa} (s.v.) defines \emph{ādāna} as, ``the three seasons of spring, summer and the rains, during which the sun takes strength from living beings''  (\emph{vasantagrīṣmavarṣās traya ṛtavaḥ ādānakālaḥ asmin sūryaḥ prāṇināṃ balam ādatte}).}
	%
	one who desires a long life should adopt at that time % tatra: may not need to be translated here?
	a regime that is more heating and moisturising.

	% ādāna: P.V. Sharma translates raukṣyam ādānajaṃ in Caraka (1.6.19cd, below) "roughness due to (beginning of) ādāna. Big help. 
	% According to the Āyurevedīyakośa (s.v.): vasaṃtagrīṣmavarṣās traya ṛtavaḥ ādānakālaḥ asmin sūryaḥ prāṇināṃ balam ādatte. So, it is the three seasons of Spring, Summer and Rains, during which the sun takes strength from living beings. 
	% Āyurvedīyakośa also says: kālādānaḥ ādadāti kṣapayati pṛthivyāḥ saumyāṃśaṃ prāṇināṃ ca balam ity ādānam
	% Cf. Caraka
	% 1.6.19ab	| hemantaśiśirau tulyau śiśire 'lpaṃ viśeṣaṇam |
	% 1.6.19cd	| raukṣyam ādānajaṃ śītaṃ meghamārutavarṣajam ||
	% 1.6.20ab	| tasmād dhaimantikaḥ sarvaḥ śiśire vidhiriṣyate |
	% 1.6.20cd	| nivātam uṣṇaṃ tv adhikaṃ śiśire gṛham āśrayet ||
	%
	% Cf. AS.Sū.4.20	
	% raukṣyaṃ ca+ādānajaṃ tasmāt kāryaḥ pūrvo 'dhikaṃ vidhiḥ | 
	% vasante dakṣiṇo vāyurātāmrakiraṇo raviḥ ||
	% Srikantha Murthy translates it as the Ādāna kāla. 
	%
	% In summarising chapter 6 of Caraka Sūtrasthāna, Meulenbeld writes:
	% The subjects are: the importance of seasonal regimen (6.3); the year is divided into six seasons; the period in which the sun courses northwards (udagayana), known as ādāna, consists of the seasons beginning with śiśira (the cool season) and ending with grīṣma (summer); the period in which the sun courses southwards, called visarga, consists of the seasons beginning with varṣāḥ (the rainy season) and ending with hemanta (winter) (6.4); the main characteristics of visarga and ādāna (6.5-8); ...
	% In footnote 119 on the main characteristics of visarga and ādāna, Meulenbeld says: During visarga the rays of the moon replenish the world, thus restoring it after a period of ādāna by the sun. 
	% When commenting on the Aṣṭāṅgasaṅgraha, Sūtrasthāna, ch. 4, Meulenbeld says: The first half of the year is called udagayana or ādāna, the second half dakṣiṇāyana or visarga. The āgneya character of ādāna and the saumya character of visarga are described (in prose, concluded by a verse) (4.5–7).

	\item[35] % not in the vulgate
	At this time (\emph{tatra}), one should avoid mixed barley grout drinks,%
	\footnote{The term \emph{mantha}, which literally means ``stirred,'' is defined in Monier-Williams dictionary (1899: s.v.) as '`mixed beverage, usually parched barley-meal stirred round in milk.'' In two instances, Ḍalhaṇa \citep[811 and 779]{vulgate} glosses \emph{mantha} as barley groats prepared with water and ghee (\emph{manthān jalaghṛtāktasaktūn}), and barley grouts covered with cold water along with ghee (\emph{manthaḥ saghṛtaśītalajalaplutāḥ saktavaḥ}).} %
	%
	frosts, windy weather, dry food, wearing dirty clothes, and bathing. %
	%
	Thus, [winter]

	\subsubsection{Spring}% vasanta
	% change numbering 

	\item[36--48] % not in the vulgate
	While the rays of the blazing sun % dīptabhāskararaśmiṣu
	are melting % nirhāra = nirhāraṇa (MW): removing
	the layers of frost, % tuṣārapaṭa
	and the days are lengthening and nights diminishing, % dineṣu jṛmbhamāṇeṣu hīymānāsu rātriṣu
	%
	when the wondrous  % vicitrāsu: variegated, diverse, but here maybe wondrous given the metaphors that follow
	rows of forests and groves % vanopavanarājiṣu
	are displaying [themselves], % darśayatsu PAP
	with colours variegated % śabalaiḥ
	by the loud laughter of flowers, % puṣpāṭṭahāsa ?
	radiant with trembling young shoots, % calat-kisalaya-ujjvalaiḥ
	and sounds % niḥsvanaiḥ
	resonating with cuckoos % kokilodghuṣṭa
	and the humming of intoxicated bees, % samadabhramarodgīta
	%
	phlegm becomes aggravated % kaphaḥ prakupyati
	and will produce % sisṛkṣati -desiderative? Whitney (948b, 1040a) says the desiderative is sometimes found where the future tense is expected.
	the diseases as described. % yathoddiṣṭān ... gadān

	\item[39] % Vulgate 6.62.32
	Cold phlegm that has accumulated in people during winter because of the cold is dissolved % viṣyaṇṇaṃ (vi-ṣyand, literally,  to dissolve, melt)
	in spring due to the heat and causes various diseases.

	\item[40]
	Because of this, one should avoid water, sweet and oily food, sleep during the day, and heavy foods, and one should also undergo treatments, such as emesis.

	\item[41--42]

	And at the end of the cold season one should eat sixty-day rice, aged rice, % ṣaṣṭikān anavāñ chālīn: could be read as sixty-day, aged rice?
	wild rice, mung beans and koda millet, as well as soup prepared properly % yuktibhiḥ - I assume this qualifies yūṣaṃ rather than adyāt.
	% include footnote on the proper method to make soup?
	with the broth of birds, such as quail, % is this lāva or lāba. In Roots of Ayurveda, lāva is bustard-quail and lāba is bush-quail.
	snake gourd, neem, rattan shoots % Ḍalhaṇa: vetrakarīraṃ vetrāgraṃ
	and Chirata% tikta: appears to be an ingredient in this list (rather than a general reference to bitter things). 
	\footnote{Chunekar and Singh (1999: 180) state that \emph{tikta} and \emph{tiktaka} (the reading of the vulgate) are synonyms for \emph{kirātatikta}, which they identify as Swertia chirata Buch.-Ham (1999: 98). Meulenbeld (2007: 31 n. 104) opined that Swertia chirata Buch.-Ham was probably a mistake for Swertia chirata (Wall.) C.B.Clarke.}
	% Meulenbeld: The Sarasvatīnighaṇṭu – Review article. Electronic Journal of Indian Medicine Volume I (2007), 21–44
	and other ingredients. % cānyair... it seems this should be read on its own?
	One should consume liquors of grape (\emph{madhu}), sugar (\emph{āsava}) and herbs (\emph{ariṣṭa}), as well as mead and sweet wine.%
	\footnote{For information on the fermented drinks called \emph{madhu}, \emph{āsava} and \emph{ariṣṭa}, see McHugh 2021: 47–49, 53. \label{liquors}}

	\item[43]
	In the Spring season,% kusumāgame: when spring flowers have arrived?
	one should adopt exercise, collyrium, smoking, and bitter medicated gargling, % kavaḍagraham: described in Suśruta 4.40.58. Should we take tīkṣṇa separately here. Ḍalhaṇa doesnt comment and PV Sharma translates tīkṣṇañ ca kavaḍagraham as 'irritant gargling.'
	and perform all these measures with warm water.

	\item[44]
	In Spring, a beneficial diet consists primarily of barley, mung beans and honey. In this season, the doubling of all exercise is recommended.
	% emendation of dvyadhikaś: is it worth a footnote? The dictionary suggests that vyadhika is often wrong for hy adhika, which is also possible here. 

	\item[45]
	Being careful, one should eliminate phlegm that has accumulated during winter with errhines, % Ḍalhaṇa: śirovirekaiḥ nasyaiḥ
	emetics, decoctions, and medicated gargling.
	% could kaṣāyaiḥ kavaḍagrahaiḥ means astringent medicated gargling?

	\item[46--47]
	At the time of Spring flowers, % sumanāḥ kusumāgame
	one who is cheerful % sumanāh
	and clean, % śuci... or does this have a more moral sense (nnocent, honest, virtuous)? 
	wearing white clothes and adorned with sandalwood and aloe-wood, should enjoy
	as much as one wants, % yāvad utsāhaṃ - lit: as far as one's resolve.
	in beautiful groves, % kānaneṣu vicitreṣu
	a woman whose breasts, thighs and buttocks are voluptuous, % emending to jaghanāṃ (K has jaṃghanāṃ and H has jaṃghānā). The emendation makes K's reading metrical (that is, the pāda becomes a na-vipulā). Also, Cf: Suśruta 6.47.63: mlānaṃ pratāntamanasaṃ manaso 'nukūlāḥ pīnastanorujaghanā haricandanāṅgyaḥ. Anyway, it has to jaghana rather than jaṃghā, as full shanks on a woman are not sexy.   
	who has beauty and youth, and is adorned with every ornament.

	\item[48]
	Then, in the evening % pradoṣaṃ as an indeclinable 
	with good friends, % satsuhṛjjanaḥ 
	one may enjoy % saṃsevet
	the light of the moon's rays, % candrakarālokaṃ: emendation: K reads candrakarālokāṃ. K's reading suggests he is supposed to enjoy a woman illuminated by the moonlight (but would he do this with good firends satsuhṛjjanaḥ?) Alternatively, if candrakarālokaṃ was correct, did the author intend us to read it with pradoṣaṃ (an evening illuminated by the moon's rays)?  
	or% athavā
	embrace one's beloved and sleep inside the house. % gṛhodare: in the interior of the house.

	% section marker in K: is this the start of summer?
	% has an iti dropped out?
	\subsubsection{Summer}% provisional ???

	\item[49-50]
	In summer, % grīṣme
	the sun takes out % ādadāna perfect middle or present middle participle of ā-da...seizing, grasping, catching hold of...
	the moisture of the world %jagatsnehan
	with its rays. % karair
	Accompanied by the sound of dry leaves stirred up % samuddhūta
	by whirlwinds, % vātacakra
	it quickly causes dryness % raukṣyaṃ sañjanayaty āśu
	in the earth, which bears a wretched, dry appearance, red with dust and embers, %
	like a woman whose bloom of youth has been ruined. % yauvanotsavanibhargnavanitāyām: disordered compound (I would expect nibhargna first). 
	Then, people dry out. % tataḥ śuṣyanti dehinaḥ

	\item[51]
	Therefore, in the hot season, one should especially avoid that which is drying and desiccating: % śoṣaṇaṃ rūkṣaṇañ ca yat
	exercise, the heat, exertion, overheated sex, % maithunaṃ paridāhi? hot sex, steamy sex? sweaty sex, vigorous sex, strenous sex? 
	and the tastes that increase the body's fire. % should I include a footnote on what these are?

	\item[52]
	In the morning, a man should drink ground meals  % saktu plural. Roots of Ayurveda is barley-meal but does the plural imply other grains as well? Ayurvedīyakośa says barley... but also yavataṇḍulalājādicūrṇam
	mixed into mango juice, % sahakāra (K,N have sahākāra, and I've corrected the spelling as sahākāra is not attested). Is "mixed with" an overtranslation of anvita here?
	sprinkled with % akta? smeared over, diffused, bedaubed...
	ghee, candied sugar, % khaṇḍa, Apte also says khaṇḍa can be a type of sugarcane
	and jaggery, % guḍa: PV Sharma's book on drugs. Roots of Ayurveda treacle?
	together with buttermilk, % Is the plural (takraiḥ) important here?
	followed by cool water. % śītatoyottaraṃ

	\item[53]
	In the hot season, % nidāghe
	one should eat preparations % vikṛtīḥ?
	of barley and wheat, and also all sorts of rice, % api suggests all...
	meats of carnivorous and marsh-dwelling animals % prasahānūpamāṃsāni
	and other aphrodisiacs, % vṛṣyāṇy anyāni ... ca
	which are prepared in various ways % prakārair vividhair 
	by a member of one's own household. % svakuṭumbinā

	\item[54]
	One should nourish % tarpayet
	desiccated bodies % śoṣitāṃs tanūn % withered, emaciated...
	with various, properly made, % vividhai... susaṃkṛtaiḥ
	well-cooled, % suśītaiḥ very cold?
	oily and sweet % suśītaiḥ snigdhamadhurais
	\emph{ṣāḍava} preparations, %
	\footnote{On \emph{ṣāḍava} preparations, see Meulenbeld 1974: 512. It is basically a condiment or confectionery made of various sweet, acid and saline ingredients.}
	% Meulenbeld Mādhavanidāna
	seasonings, % rāgair ?
	and fermented jaggery drinks. % gauḍikaiḥ | Ḍalhaṇa: gauḍikaṃ guḍakṛtaṃ madyam

	\item[55]
	In summer, food that is wholesome, cold, sweet, and liquefied, % drava
	together with sugar-sweetened milk, % saśarkareṇa payasā
	is also recommended at night.

	\item[56--57]
	During the hot season, one should frequent lakes, ponds, % vāpyaḥ: could be more specific: dams or reservoirs
	rivers, and pleasant forests, as well as excellent sandalwood trees, % candanāni parārdhyāni
	garlands with lotuses and water lilies, % I've read against K and N by adopted srajaḥ. Both K and N have sakamalotpalāḥ, feminine acc, which needs srajaḥ to make sense. 
	the breeze of palm-leaf fans,  % \skt{tāla-vṛnta} n. a palm-leaf used as a fan, fan
	necklaces, % hārāṃs : cf vulgate.
	cool houses, % śītagṛhāṇi
	and very light clothes.

	\item[58]
	During the day, one should lie down as one pleases % yathākāmaṃ
	in a cool house adorned with lotus and water-lily petals, % padmotpaladalāyute
	while being pleasantly touched by breezes. % spṛśyamāno 'nilaiḥ sukham: by the a breeze might sound better in English. 

	\item[59--60]
	At night, surrounded by friends, % suhṛdbhir abhisamvṛtaḥ
	one should sit in a charming % hṛdye is charming too weaK? cherished?
	palace terrace, % harmyatale : palace roof or  upper room MW
	decorated with scattered heaps % prakara
	of flowers, along with variegated lotus blossoms, % jalajaih kusumaiś citraiś. MW \skt{-kusuma} n. `water-flower', lotus, 
	and covered by the moonlight.
	One should be fanned and touched with soft fans, % vyajanair vījyamānaś ca spṛśyamānaś ca komalaiḥ
	while being attended on by women % strīṇāṃ
	with their breasts, hands and cool pearl necklaces.

	\item[61] % there is a section marker in K and H at the end of this verse.
	When bodies are dry and dessicated or when it is summer, % rūkṣocchuṣkeṣu deheṣu grīṣme vā... The pl. deheṣu is odd.
	a man whose breath has the floral odour of impending death  % puṣpito % MW: exhaling an odour indicative of approaching death Caraka. Āyurvedīyakośa: iṣṭāniṣṭapuṣpagandhasādṛśyanirūpitariṣṭalakṣanalakṣitaḥ (ci. 2.8-10)
	and is weak % abalī
	should rest in a beautiful bed, with a blanket moistened by sandal. % puṣpitovalī

	% K, H... section marker here but no iti.

	\subsubsection{Monsoon}% ??

	\item[62--65] % 64.46.ADD8
	In the monsoon, % prāvṛṣi : in the rainy season?
	when the sky % ambara?
	resonates loudly % ativinādite
	with clouds whose roar is deep and great,  % pravṛddha-gambhīra-ravair
	%
	causing the \gls{kandala}, % MW: N. of a plant with white flowers (which appear very plentifully and all at once in the rainy season) Suśr. Ritus. Bālar. Possibly related to Kandalī
	\glspl{śyāma}, % śyāmā RoA or priyaṅgu. MW: śyāma: N. of various plants (fragrant grass; thorn-apple; Artemisia Indica; Careya Arborea &c.) L
	\gls{kadamba}, % RoA: Neolamarckia cadamba (Roxb.) Bosser: kadamba,śyāma-kadamba: unattested?
	\gls{kuṭaja} % RoA kuṭaja
	and \gls{arjuna} to burst forth; % ārjuna: RoA: arjun, Terminalia arjuna, Bedd.:
	%
	and filled % āpūrite
	throughout its ten regions  % daśāśā ten regions. See {daśa-diś} f. sg. the 10 regions
	with the perfume of Jasmine flowers;% mālatīkusumāmoda: RoA: mālatī, jasmine,
	%
	when the surface of the earth % {vasdudhā-tala} n. the surface of the earth 
	is covered % prakīrṇṇa:
	with \glspl{surendragopa} % velvet-mites % surendragopa. 
	and frogs, % maṇḍūka
	and adorned by % bhūṣite
	music of black bees % dvirephakṛtasaṅgīta
	and groves of lotuses;
	%
	when the festival begins with the cries of the foremost peacocks; % too free: literally: when the festival produced by the commencement of cries of the best leaders of the peacocks. 
	and the sun reaches the uppermost limit of its weakness, % daurbalyakāṣṭhā?
	%
	wind is aggravated by the cold, wind, clouds, and rain.

	\item[66]
	Without delay % avilambena
	and by the proper method, % vidhinā
	one learned in diseases % vyādhikovidaḥ: An expert in pathology?
	should eliminate % nihanyād
	the wind accumulated during the hot season % nidāghopacitaṃ
	that is being aggravated. % prakupyantaṃ

	\item[67]
	One should take % caryāḥ
	tepid % sukhoṣṇāś 
	broths % rasaḥ ??
	and various oils, % tailāni vividhāni ca
	% Cf:   6.64.47 ab| tāpātyaye hitā nityaṃ rasā ye guravas trayaḥ||
	%       6.64.47 cd| payo māṃsarasāḥ koṣṇās tailāni ca ghṛtāni ca|
	%       6.64.48 ab| bṛṃhaṇaṃ cāpi yat kiñcid abhiṣyandi tathaiva ca||
	% At the end of the summer season, the three heavy tastes are always beneficial. And milk, tepid meat soups, oils, and ghees, as well as anything that is nourishing and causes defluxion.
	as well as anything fattening % bṛṃhaṇa
	and causing defluxion. % abhiṣyandi % MW running eyes (Suśruta)

	\item[68]
	One should avoid fresh water, % navāmbu 
	dry and cold food, and ground meals.% saktu
	One should enjoy barley, sixty-day rice, % is ṣaṣṭikaśālīṃś sixty-day rice and rice?
	wheat, and small grains, % anu = aṇu ? (4) name of very small grain such as sarṣapa, cīnaka &c., 
	as well as a very soft bed in a draught-free room within the palace.

	\item[69]
	Rain water % toyam āntarīkṣaṃ
	that is filled with % samāplutan: inundated by?
	the faeces, urine, saliva, % lālā
	expectoration, % niṣṭhīvana  (is the difference between lālā and niṣṭhīvana to do with the former being saliva and the latter phlegm from the throat and chest?)
	and so forth of poisonous creatures, then becomes like poison. %

	\item[70]
	Rainwater is contaminated % dūṣitam (supplying āntarīkṣaṃ toyaṃ)
	by monsoon wind % prāvṛṣeṇyena, vāyunā 
	that is poisonous. % viṣajuṣṭena 
	Therefore, % hi
	during that season one should avoid such water for any purpose.

	\item[71]
	In the monsoon season, % prāvṛṣi
	one should judiciously % yuktitaḥ ... judiciously
	consume aged % jīrṇān
	herb, sugar and spicy liquors
	%\skt{maireya}    m. n. a kind of intoxicating drink (accord. to Suśr. Sch. a combination of \skt{surā} and \skt{āsava}) MBh. R. &c.
	% upadaṃśa m. Anything which excites thirst or appetite, a relish, condiment 
	with condiments.\footnote{On herb (\emph{ariṣṭa}) and sugar (\emph{āsava}) liquors, see footnote \ref{liquors}. \emph{Maireya} (``spicy liquor'') is similar to \emph{āsava}, but made from various other sugars and spices. \citet[50--53]{mchu-2021} describes it as an ancient, spiced, variable-sugar alcoholic drink.} % sopadaṃśāṃś
	One should avoid them at night during the monsoon.

	\item[72]
	One should alleviate aggravated wind with purging % nirūhair
	enemas as well as other wind-removing treatments. One should also follow the monsoon precepts.

	\item[73]
	However, rain % varṣī
	clouds % parjanya - collective singular?
	are more abundant % bhūya
	for people south of the Ganges.
	There, two rainy seasons called Prāvṛt and Varṣa are known to them. % not sure why the genitive plural is used if the referent is jana?

	\item[74]
	In its northern region abounding in snow and ice, % himavatdravya: literally: snowy substances
	it is much % bhūyaḥ
	colder there during early % hemanta 
	and late winter.% śiśira

	\item[75]
	A man who lives by the prescribed regime in each season, %
	never % na kadā cana. But K does not have the negative. K is wrong. But is it?
	incurs the terrible diseases caused by the seasons.

	\subsubsection{The Twelve Diets}

	\item[76] % 6.64.56
	Next, % ata ūrdhvaṃ
	we shall describe the twelve classifications of diet. % aśana 
	These are the following: % tad yathā
	\begin{itemize}
		\item cooling, %
		\item heating,
		\item oily,
		\item oil-free,\footnote{
			      On the meaning of \emph{rūkṣa} in the context of diet, see  Wujastyk 1998 263 n. 39.
		      }
		\item liquid,
		\item desiccating, % śuṣka? 1-6 vs rūkṣa % dry 
		\item once a day, % ekakālika
		\item twice a day, % dvikālika
		\item combined with medicines, % auṣadhayukta 
		\item reduced-quantity, % hīnamātra
		      and \item those with the aim of pacifying and % doṣapraśamana 
		\item increasing the humours. % pravṛddhi - nourishing? PV Sharma says: maintenance of body but the vulgate has vṛtty here.
		      % Not sure about arthān... Sharma does not translate it. Was it intended to mean simply the topics or matters of ... and should it be construed with the whole list? or maybe as I have read it (pravṛddhyartha).
	\end{itemize}

	\item[77]
	One should treat with cooling foods % śītair annair	
	those suffering from thirst, derangement, % mada?
	and inflammation, and % vidāha 
	those afflicted by \se{raktapitta}{choleric blood} % raktapitta. RoA
	and poison,
	as well as fainting in women and emaciated people. % kṣīṇān is well attested elsewhere. 

	\item[78] % xml:id 6.64.58
	One should treat with heating foods % uṣṇair annair 
	those burdened with% āviṣṭa
	phlegm, wind and disease;
	those who have been purged; % aviktān
	those consuming oily drinks; % snehapāyinaḥ
	and men whose bodies are not moist. %aklinnakāyāṃś ca narān
	% the question to ask here whether narān should be read with all the adjectives and, if so, whether these verses are being specific about gender given strīṣu in the previous verse. PV Sharma ignores narān.

	\item[79] % xml:id 6.64.59
	One should treat with oily foods % snigdhair annair 
	men affected by wind-diseases, % vātikān
	those whose bodies are dry,
	who have been weakened by sex, % vyavāyopahatāṃs 
	and also those who exercise.

	\item[80] % xml:id 6.64.60
	One should treat with oil-free foods % rūkṣair annair
	those who have become fat, % medasābhiparītāṃs : literally: those filled with fat.
	who are corpulent, % sthūlān % stocky?
	who are suffering urinary disease,
	and whose bodies have been overwhelmed with phlegm. % kaphābhipannadehāṃś 

	\item[81] % xml:id 6.64.61
	With liquid foods, one should treat those who have dry bodies, are suffering from thirst, and are also weak. %
	With drying foods, those whose bodies are moist, who are wounded and those with urinary disease.

	\item[82]
	Food should be given once a day to increase weak digestive fire. %
	And % atha
	it is recommended % pūjitaḥ
	twice a day % dviḥkālam : the visarga is unconventional but maybe okay as dvis + kāla
	to equilibrate the fire. % samāgnaye balance, harmonise

	\item[83] % xml:id 6.64.63
	For a person who dislikes medicine, one should give food mixed with the medicine. % obviously the word to be supplied for food here was neuter, such as bhojana.
	A diet reduced in quantity is recommended for weak digestive fire and for one who has a disease.

	\item[84] % xml:id 6.64.64
	Food that is given according to the season is regarded as pacifying the humours. %
	After that, any diet is recommended for healthy people.

	\subsubsection{The Times for Administering Medicine}

	\item[85] % 64.67
	Next, we will explain the ten times for administering medicine. In this regard, they are:
	\begin{itemize}
		\item without having eaten,
		\item before eating,
		\item after eating,
		\item in the middle of eating,
		\item between meals,
		\item together with a meal,
		\item before and after a meal, % sāmudgaṃ MW (from samudga - casket) enclosing the food. 
		\item taking it repeatedly, % muhur muhur
		\item as a mouthful, % grāsa
		\item and between mouthfuls. % grāsāntaraṃ
	\end{itemize}

	\item[86]
	Among these, % tatra
	the one called "without having eaten" is that in which % yaḥ must be a mistake (as nothing in this sentence is masc.) so I have adopted the vulgate's yat. It is going to take me a while to come down from the thrill of having improved the Nepalese text. 
	medicine alone is taken.

	\item[87]
	Medicine without food becomes more potent and it quickly thus cures a disease, without doubt. %
	However, when children, the elderly, women, and weak persons have taken it, % tad ... pitvā
	they become extremely fatigued and lose their strength.

	\item[88]
	The one called "before eating" is that in which medicine is taken before eating. %
	Food is eaten afterwards. % paścād bhujyate tat - the referent of tat is not clear grammatically but surely must be food. 

	\item[89] % 64.69
	Medicine served before eating % prāgbhaktasevitam: in the Nepalese version, we have to supply medicine. The vulgate has athauṣadham etad eva.
	is quickly digested, % śīghraṃ vipākam upayāti
	does not impair one's strength, % balaṃ na hiṃsyād
	and, when embedded in food, % annāvṛtan ... ca
	is not repeatedly % muhur ?
	regurgitated from the mouth.
	It gives strength below,% adhobalam? 
	\footnote{Although not attested elsewhere in medical literature, the term \emph{adhobala} (literally, ``lower strength'') appears to mean strength in the lower part of the body. This follows a pattern seen in the next verses where medicine taken after a meal affects the upper body (\emph{ūrdhvakāya}) and medicine taken in the middle of a meal affects the middle of the body (\emph{dehamadhya}).}
	% I am tempted to adopt athauṣadham etad eva, as it states the subject clearly. And adhobalam could be a corruption of athoṣadhaṃ (on orthographic grounds) and ādadhāti could be a dittographic error (as the next pāda begins with dadyāt). However, adhobalam ādadhāti might be a truncated expression parallel to ūrdhvakāye ... balan dadāti.
	% In the third quarter, we have adopted the vulgate reading of <foreign>athauṣadham etad eva</foreign>, which states the subject of the verse clearly. The meaning of the Nepalese reading, <foreign>adhobalam ādadhāti</foreign> (literally, ``it gives lower strength'')  is not clear. The compound <foreign>adhobala</foreign> is not attested, as far as we are aware, in Sanskrit literature.
	%
	and should be given to the elderly, children, the timid, those who are emaciated, and women.

	\item[90]
	The one called "after eating" is that in which medicine is taken when one has eaten. %

	\item[91]
	Having taken medicine % yad upayujya
	when food has been consumed, % pītaṃ ... annam 
	it % tad
	cures various diseases and gives strength in the upper body. %

	\item[92]
	The one called "in the middle of eating" is that in which medicine is taken in the middle of eating. %

	\item[93]
	As for medicine taken in the middle of eating % madhye tu pītam 
	it cures % apahanty
	diseases %  supply tān with ye rogāḥ
	whose nature is not to spread % avisāribhāvā
	and which %
	affect the middle of the body. % dehamadhyam abhibhūya

	\item[94]
	The one called "between meals" is that in which medicine is taken in between two meals, the morning and evening.%
	\footnote{Ḍalhaṇa glosses ``earlier and later'' as morning and evening meals (\emph{pūrvāparayor iti prātaḥsāyantanayor bhojanayor ity arthaḥ}). }

	\item[95]
	And medicine % tat... ca
	taken between meals %antarabhaktakan
	is agreeable, % hṛdya... 
	strengthens the mind, stimulates digestion, and is always wholesome.

	\item[96]
	The one called "together with a meal" is that which is administered in regard to medicines.%
	\footnote{This statement was rewritten more clearly in the vulgate: The one called "together with a meal" is that in which [the medicine is taken] with food (\emph{sabhaktaṃ nāma yat saha bhaktena}).}

	\item[97]
	Medicine taken together with a meal % sabhaktam ... tat
	is always wholesome for women, % abalā
	the weak, % abala
	and all those who dislike medicine, % dveṣiṇām api tathā
	as well as children and the elderly.% śiśuvṛddhayoś ca 
	\footnote{Ḍalhaṇa glosses \emph{abalābalayoḥ} as ``for women and the weak'' (\emph{strīdurbalayoḥ}), and \emph{dveṣiṇām} as ``for those who dislike medicine'' (\emph{bheṣaja-dveṣiṇām}).}

	\item[98]
	The one called "before and after a meal" is that in which medicine is taken at the beginning and end of a meal.

	\item[99]
	When a humour has spread two ways,% pravisṛte vulgate. K's reading is unmetrical, H's reading is not attested
	\footnote{Ḍalhaṇa clarifies that ``two ways'' (\emph{dvidhā}) means up and down (\emph{dvidhā pravisṛte ūrdhvādhaḥprasṛte ity arthaḥ}).}
	that called "before and after a meal", % samudgasaṃjñam
	which is at the beginning and end of a meal, % ādyantayor yad aśanasya
	is administered. % niṣevyate: I assume the final tu is pādapūrṇa

	\item[100]
	The one called "taking it repeatedly" is that in which medicine is taken repeatedly, either with or without food.

	\item[101]
	When repeated % muhur muhur 
	wheezing % śvāsa
	cough, hiccups and vomiting % kāse hikkāvamīṣu
	arise excessively,  % atiprasṛte
	this regime %etad
	should be administered.

	\item[102]
	The one called "a mouthful" is that in which medicine is taken mixed as a lump of food.%
	\footnote{Ḍalhaṇa adds that as seen in other treatises, it is understood as mouthful by mouthful (\emph{tantrāntaradarśanāt grāsaṃ grāsam iti boddhavyam}).}

	\item[103]
	One should endeavour to administer in mouthfuls powered medicine that is a digestive stimulant % %
	for those with weak digestive fire,
	and also as an aphrodisiac.

	\item[104]
	The one called "between mouthfuls" is that in which medicine is taken between lumps of food.

	\item[105]
	Between mouthfuls, %grāsāntareṣu
	one should give out % visṛjet : basically means administer/give but I'm looking for a different word in English to reflect the different word in Sanskrit
	emetic smokes % vamanīyadhūmāñ
	and, in cases of wheezing and the like, %chvāsādiṣu
	electuaries % lehān syrups
	whose properties are known and observed. % prathitadṛṣṭaguṇāṃś

	\item[106]
	These are the ten times for administering medicine. And on this topic, there is the following:

	\item[107]
	When faeces and urine have been discharged % visṛṣṭe viṇmūtre
	and the body is very light with the senses calm; % viśadakaraṇair dehe ca sulaghau
	when eructation is clear, the heart pure, and flatus is moving through;
	% ḍalhaṇa: āte saratīti adhaḥ pravahati vāte ity arthaḥ 
	and when there is desire for food, % tathānnasraddhāyāṃ
	fatigue has gone, % Ḍ: klamaparigame iti anāyāsaśramarahite ityarthaḥ
	and the abdomen is relaxed,
	the physician should give food. % Ḍ: āhāro dātavyo bhavati bhiṣajā vaidyena,
	This is a wholesome observance. % prakriyā : procedure, method, observance, prescription
	% I do not recognise the metre of this verse, but it appears to scan, except for the end of the fourth pāda. is it a faulty Śārdūlavikrīḍita
	% lggggglllllggglllg    {m: 27}    [18: ymnsts]
	% lggggglllllgggllll    {m: 26}    [18: ymnstn]
	% lggggglllllggglllg    {m: 27}    [18: ymnsts]
	% lggggglllllgggglgl    {m: 28}    [18: ymnsmj]

	\item[108]
	The Lord Dhanvantari said.
\end{translation}
